\subsection{Gleichrichterschaltung}

\Aufgabe{
  \label{Aufg28}
  Hier eine Gleichrichterschaltung mit pn-Silizium-Dioden.
  Die Spannungsquelle $V_1$ liefert sinusförmige $50 \, \mathrm{Hz}$ Wechselspannung mit einer Amplitude von $5\, \mathrm{V}$.
  \begin{figure}[H]
    \centering
    \input{Tikz/src/FigKlausurGleichrichter}\alttext{Das Bild zeigt ein elektrisches Schaltbild mit einer Wechselspannungsquelle \(u_\mathrm{E}(t)\), die an zwei antiparallel geschaltete Dioden angeschlossen ist. Jede Diode ist parallel zu einem Kondensator geschaltet, wobei der obere Kondensator \(C_1\) mit dem Widerstand \(R_1\) in Reihe und der untere Kondensator \(C_2\) mit dem Widerstand \(R_2\) in Reihe verbunden ist. Beide Reihenstränge sind parallel zueinander geschaltet und enden in einem gemeinsamen Widerstand \(R_3\). Die Spannungsquelle ist geerdet.}
    \label{fig:FigKlausurGleichrichter}
  \end{figure}
  \begin{itemize}
    \item[a)]
          Welche Gleichspannungen treten über $R_1$,$R_2$ und $R_3$ auf?
    \item[b)]
          Welche Funktionen haben die Kondensatoren $C_1$ und $C_2$?
    \item[c)]
          Was ist der besondere Vorteil dieser Gleichrichterschaltung, z.B. für Anwendungen mit Operationsverstärkern?
    \item[d)]
          Skizzieren Sie eine andere Spannungs-Verdopplungs-Schaltung, bei der die maximale positive Gleichsspannung (Ausgangsspannung) gegenüber Massepotential der speisenden Wechselspannung $V_1$ zur Verfügung steht.
  \end{itemize}
}

\Loesung{
  \begin{itemize}
    \item[a)]

          $U_\mathrm{R1} = 4,3 \, \mathrm{V} \Rightarrow$ da $0,7 \,\mathrm{V}$ über $D_1$ abfallen \\
          \\
          $U_\mathrm{R2} = -4,3 \, \mathrm{V} \Rightarrow$ da gleiches für die negative Halbwelle \\
          \\
          $U_\mathrm{R3} = 8,6 \,\mathrm{V} \Rightarrow$ da $R_3$ beide Spannungswerte aufnimmt \\


    \item[b)] Die Kondensatoren $C_1$ und $C_2$ haben einerseits die Funktionen die Signale zu glätten, sowohl von der negativen als auch der positiven Halbwelle.

    \item[c)] Der Vorteil dieser Gleichrichtung ist, das sowohl positive als auch negative Spannungen abgegriffen werden können.

    \item[d)] Schaltung
          \begin{figure}[H]
            \centering
            \input{Tikz/src/LsgKlausurGleichrichter}
            \alttext{Die Abbildung zeigt eine Schaltung aus zwei Dioden und zwei Kondensatoren, die zwischen einer Eingangs- und einer Ausgangsklemme angeordnet sind.
            Links ist die Eingangsspannung \(u_\mathrm{e}\) eingezeichnet, die zwischen dem oberen linken Knoten und dem unteren Bezugsknoten gemessen wird. Der untere Leiterzug stellt den gemeinsamen Massebezug dar und verbindet alle unteren Knoten miteinander.
            Vom oberen linken Eingang führt zunächst der Kondensator \(C_1\) in Serie zu einem Zwischenknoten. Von dort ist die Diode \(D_2\) in Durchlassrichtung zum rechten oberen Knoten geschaltet. Dieser rechte obere Knoten bildet die Ausgangsklemme.
            Zwischen dem Zwischenknoten vo \(C_1\) und dem unteren Leiterzug ist die Diode \(D_1\) geschaltet. Weiter rechts ist zwischen dem oberen rechten Knoten und Masse der Kondensator \(C_2\) angeschlossen.
            Die Ausgangsspannung \(U_\mathrm{a}\) wird zwischen dem rechten oberen Knoten und dem unteren Bezugsknoten gemessen.}
            \label{fig:LsgKlausurGleichrichter}
          \end{figure}
  \end{itemize}
  Siehe: Abschnitt \ref{fig:BrueckengleichrichterSchaltung}
}